This dissertation represents the culmination of five years of research. Still, beyond the equations, simulations, prototypes, papers, and presentations, it stands for something far more personal. These years have been a journey of self-discovery and growth; one that has challenged me, humbled me, and ultimately shaped me into a better researcher, teacher, colleague, and person. None of this would have been possible without the many wonderful people who supported me, believed in me, and shared this journey with me along the way.

First and foremost, I want to thank my advisor and mentor, \textbf{Dr. Evgueni Filipov}. I first met you as a student in your Structural Dynamics class during my Master's, and one of the first things I remember about you has little to do with dynamics itself. You would take music suggestions from students to play during the first few minutes of class while everyone settled in. It was a gesture that captured something I would come to appreciate about you over the years: your ability to make people feel comfortable and welcome. Your Reconfigurable Structures class later sparked my interest in deployable and reconfigurable structures and led to an independent study that eventually became one of the projects and papers in this dissertation. Looking back, I could not have known how consequential those first few classes and conversations would become.

I had no publications, little experience with independent research, and only a vague understanding of the field when you gave me the opportunity to join your lab. In the uncertainty of COVID, I was also unsure whether I wanted to pursue a PhD then or later. Despite that, you believed in me and gave me a chance to find out what I was capable of. That decision truly changed the trajectory of my life. I still remember complaining that I hated coding during my early days, as you patiently encouraged me to keep trying. Considering how central computational simulations and programming eventually became to my PhD, the transformation I went through still makes me smile. You changed the way I thought about research and the kinds of problems I believed I could solve.

Perhaps the most valuable thing you gave me was the freedom to find my own way. You rarely told me exactly what to do; instead, you had a way of asking the right questions and letting me figure out the next step for myself. I often came to your office in the early days of my Ph.D., hoping you would tell me what to do next. Somewhere along the way, I began coming in wanting to discuss what I thought we should do next. You helped me develop the confidence to trust and defend my ideas while remaining open to criticism. You also taught me that doing good research is only a job half done; what completes it is communicating that research in a clear and accessible manner. I have lost count of how many times you reminded us to read Michael Alley's \textit{The Craft of Scientific Writing} and \textit{The Craft of Scientific Presentations}. I now find myself returning to those lessons whenever I write a paper, prepare a talk, or explain my work to someone outside the field.

I am equally grateful for your patience and support when research, teaching, or life in general became tough. Whether we were debugging software and trying to make sense of confusing results, or I was overwhelmed by the load of teaching for the first time ever, you never made slow progress feel like failure. Your calmness taught me to take difficult problems one at a time, while your optimism and willingness to explore new ideas made research exciting. You always cared about your students as people, not just researchers. You encouraged us to take vacations, pursue hobbies and internships, explore different careers, and make time for ourselves and our families. As an international student, I was especially grateful for the understanding and flexibility you showed around travel, visas, and all the complications that came with them. That same care was reflected in the culture you created in the lab, from celebrating everyone's milestones and organizing outings to inviting us into your home and, of course, bringing us bottles of honey from your beehives. Your mentorship helped me navigate research, teaching, publishing, and career decisions, eventually supporting my path outside academia, even if you still hope I find my way back someday. By the end of this PhD, our conversations felt less like those between a student seeking direction and an advisor providing it, and more like conversations between colleagues. I hope to carry your curiosity, optimism, and generosity with me throughout my career. Thank you for taking a chance on me, for giving me room to grow, and for caring about the person behind this dissertation. Other members of our lab and I have been incredibly fortunate to have you as our advisor, and I will always be proud to have been your student.

I would also like to thank my committee members and collaborators, \textbf{Dr. Jeffrey Scruggs}, \textbf{Dr. Kevin Maki}, and \textbf{Dr. Lorenzo Valdevit}. Each of you brought a unique perspective to my work and encouraged me to think beyond the boundaries of my own expertise. Your questions often pushed me to examine my assumptions carefully, consider problems from a different point of view, and communicate the broader significance of my work clearly. I am especially grateful for the thoughtful feedback you provided on this dissertation and on the papers and research ideas I developed throughout my PhD. Your suggestions and constructive feedback have made this dissertation its strongest self.

To the members of the \textit{Reconfigurable, Architected, and Regenerative (RARe) Structures Lab}, formerly known as the \textit{Deployable and Reconfigurable Structures Lab (DRSL)}: Thank you for creating a community where I could learn, grow, and feel at home. \textbf{Yi}, \textbf{Wo}, \textbf{Steven}, and \textbf{Maria}, it has been an honor to follow in your footsteps. I learned so much simply from watching how you approached research and from the knowledge and experience you shared with me. The especially high standards you set gave me something to aspire to throughout my Ph.D. and your feedback and advice continued to help me long after our time in the lab overlapped. \textbf{Anan}, \textbf{Anvay}, \textbf{Jacob}, \textbf{Raj}, \textbf{Wayne}, \textbf{Vaibhavi}, \textbf{Giwan}, and \textbf{Kwanwoo}: Thank you for the conversations, the questions, the feedback, and the countless times you allowed me to bounce unfinished ideas off you. Some of the most useful discussions began when I stopped by with a problem I had not yet fully understood. Thank you for listening, challenging my assumptions, and helping whenever I needed it. It has been a genuine pleasure to work alongside such thoughtful and talented people. I am grateful to call you my colleagues and friends.

Special thanks to \textbf{Daniel Hart} and \textbf{Jason Kelly} from the United States Navy's Naval Surface Warfare Center, Carderock Division (NSWCCD). Our collaboration on the curved-origami deployable hulls project took my research into an entirely new area and became one of the most exciting parts of my doctoral journey. Thank you for the many helpful discussions, technical insights, and thoughtful questions that helped connect our ideas to a meaningful engineering application.

I am also deeply grateful to \textbf{Dr. Robert Sulewski} and \textbf{Dr. Ann Jeffers}, who were the best colleagues and mentors I could have asked for while serving as a Graduate Student Instructor for the ENGR-100 and CEE-212 courses, respectively. Working with you gave me a deep appreciation for the art of teaching. You showed me that good teaching is not only about explaining technical ideas, but also about building confidence, creating curiosity, and helping students recognize what they are capable of. Thank you for trusting me, supporting me, and giving me the space to develop my own voice in the classroom. Teaching was one of the most rewarding parts of my graduate student experience, and I will always be grateful for the skills and confidence I gained while working with you.

\textbf{Mom} and \textbf{Dad}, thank you for your unconditional love and unwavering support throughout this journey, even when you did not fully understand what I was working on. You supported every decision, celebrated every small victory, and reminded me that I always had a home to return to. Pursuing this degree meant building a life far away from home, but your love and encouragement made the distance feel smaller. Everything I have accomplished is rooted in the opportunities, sacrifices, values, and confidence you gave me. I hope this moment makes you immensely proud.

\textit{``Whatever you do in this life, it's not legendary unless your friends are there to see it."} To my friends in India and in the United States, thank you for making sure I never had to go through this journey alone. You kept me grounded, gave me something to look forward to when research became overwhelming, and constantly reminded me that there was a life outside research. You were there for the occasional existential crises that seem to come with every Ph.D., the breakthroughs, and all the moments in between. \textbf{Anshul, Aditya, Riya, Sanyukta, Srikanth, and Tarun}, I am so grateful that my time in Ann Arbor began with all of you. Even after you moved, the trips, visits, and adventures we planned always gave me something to look forward to. \textbf{Nishant}, thank you for ensuring there was never a shortage of parties, food crawls, Marathi theatre, or some new plan to drag us into. \textbf{Shrey} and \textbf{Dyuti}, thank you for the countless game nights that made even ordinary weeks memorable. And \textbf{Sonal}, I am especially grateful that somewhere along the way you became one of my closest friends in Ann Arbor. From being my gym partner to being ready for almost any spontaneous plan, your friendship made Ann Arbor feel like home again. \textbf{Pratyarth}, thank you for always cheering me on and for being there through every version of me along the way since our undergrad days. After all these years, I am glad we finally get to be in the same city as this chapter ends. \textbf{Bhairavi}, even when distance and life meant we spoke less often than we would have liked, every conversation felt like we were picking up where we left off. That familiarity meant more than you probably realize during all these years away from home.

Graduate school could easily have been lonely, but because of all of you, it never was. I may have missed mentioning a few people by name, but please know that this does not diminish my gratitude for your support and presence throughout these years. Thank you for the laughter, adventures, celebrations, and companionship that made these five years about so much more than earning a Ph.D.

Last, but certainly not least, \textbf{Manvi} \textemdash{} Thank you! In more ways than one, you have had a front-row seat to the entire journey of my Ph.D. You have seen the excitement, the doubt, the exhaustion, the last-minute changes, and everything in between. Through it all, you remained my biggest cheerleader and the person I could rely on in both the highest and lowest moments. Thank you for every late night you stayed up with me, helping me perfect a presentation, revise an application, rehearse a talk, or work through whichever challenge had taken over my life that week. Thank you for listening patiently to ideas that were often far too technical, for knowing when I needed advice and when I needed someone beside me, and for believing in me during the moments when I found it difficult to believe in myself. So many of the things I am proud of over the five years of my Ph.D. carry your fingerprints in one way or another. I count myself extremely lucky to call you my partner and my best friend. So excited to see where life takes us next!

\begin{center}
    \rule{0.5in}{0.25pt}
\end{center}

The research presented in this dissertation was supported in part by the \textbf{Automotive Research Center Cooperative Agreement W56HZV-19-2-0001}, the \textbf{National Science Foundation Grant No. 1943723}, and the \textbf{Office of Naval Research Grant No. N00014-20-5-B001 \#13209604}. Any opinions, findings, conclusions, or recommendations expressed in this dissertation are those of the author and do not necessarily reflect the views of these institutions.